% ============================================================
% Milestone 1 -- Project Proposal & Literature Review 
% DIP Team Project
% ============================================================
\documentclass[11pt]{article}

\usepackage[margin=1in]{geometry}
\usepackage[T1]{fontenc}
\usepackage[utf8]{inputenc}
\usepackage{amsmath,amssymb}
\usepackage{graphicx}
\usepackage{booktabs}
\usepackage[round]{natbib}
\usepackage{hyperref}
\usepackage{enumitem}
\usepackage{xcolor}

\hypersetup{
  colorlinks=true,
  linkcolor=blue!50!black,
  citecolor=blue!50!black,
  urlcolor=blue!50!black
}
\newcommand{\todo}[1]{\textcolor{red}{\textbf{[TODO: #1]}}}

\title{\LARGE Project Proposal \& Literature Review\\[4pt]
  \large Lightweight Learned Image Compression for Laptop-Class Training}
\author{\todo{Student A} \and \todo{Student B} \and \todo{Student C}}
\date{Milestone 1\\CSE681 --- Digital Image Processing}

\begin{document}
\maketitle

\begin{abstract}
\todo{In 3--5 sentences, state the problem, motivation, planned approach, and
dataset and baseline.}
\end{abstract}

\section{Problem Statement and Motivation}
\todo{Describe the problem, project scope, motivation, and connection to a course
topic.}

\section{Related Work}
\todo{Summarize at least four references, including at least two from 2023--2026.
Explain how the work informs this project.}

\begin{table}[h]
\centering
\caption{Key references and their relevance to this project.}
\begin{tabular}{@{}p{0.36\linewidth}p{0.12\linewidth}p{0.42\linewidth}@{}}
\toprule
\textbf{Reference} & \textbf{Year} & \textbf{Relevance} \\
\midrule
\citep{balle2018scalehyperprior} & 2018 & Scale-hyperprior learned image codec. \\
\citep{he2022elic} & 2022 & Efficient learned image compression architecture. \\
\todo{Recent reference} & \todo{Year} & \todo{Relevance} \\
\todo{Recent reference} & \todo{Year} & \todo{Relevance} \\
\bottomrule
\end{tabular}
\end{table}

\section{Dataset and Baseline Plan}
\todo{Specify the DIV2K training subset and crop size, Kodak/Tecnick evaluation
sets, dataset links and licenses, and the JPEG baseline at matched bits-per-pixel.}

\section{Proposed Method}
\todo{Describe the reduced-width hyperprior codec, planned width or latent-channel
settings, and include an architecture figure or pseudocode.}

\section{Evaluation Plan}
\todo{Define how bits-per-pixel (bpp), PSNR, and MS-SSIM will be measured and how
the rate--distortion curves will be compared.}

\section{Feasibility and Compute Plan}
\todo{State the compute tier, training setup, expected experiment time, and result
of the one-training-step feasibility test.}

\subsection{Minimum Viable Deliverable}
\todo{State the smallest working experiment that produces a valid comparison with
JPEG.}

\subsection{Stretch Goal}
\todo{State what the team will attempt if the MVP is on schedule by Week 5--6.}

\section{Risks and Mitigations}
\begin{itemize}[leftmargin=*]
  \item \todo{Risk and mitigation or fallback.}
  \item \todo{Risk and mitigation or fallback.}
\end{itemize}

\section{Work Breakdown}
\begin{table}[h]
\centering
\caption{Task allocation by milestone and team member.}
\begin{tabular}{@{}p{0.16\linewidth}p{0.25\linewidth}p{0.25\linewidth}p{0.25\linewidth}@{}}
\toprule
\textbf{Milestone} & \textbf{Student A} & \textbf{Student B} & \textbf{Student C} \\
\midrule
M1 (Week 3) & \todo{Tasks} & \todo{Tasks} & \todo{Tasks} \\
M2 (Week 6) & \todo{Tasks} & \todo{Tasks} & \todo{Tasks} \\
M3 (Week 9) & \todo{Tasks} & \todo{Tasks} & \todo{Tasks} \\
\bottomrule
\end{tabular}
\end{table}

\bibliographystyle{unsrtnat}
\bibliography{references}

\end{document}
